Le condensateur ayant acquis sa charge maximale $q_{\max }$. Lorsqu'on bascule l'interrupteur en position 2 (Figure 1) à un instant $t_{0}=0$ choisi comme nouvelle origine de temps, le condensateur se décharge. La figure 3 représente l'évolution temporelle de la tension $u_{C}(t)$.\\
Donnée : $r=8,6 \Omega$\\
2.1. Établir l'équation différentielle vérifiée par la tension $u_{C}$.\\
2.2. Nommer le régime d'oscillation obtenu.\\
2.3. Déterminer graphiquement la valeur de la pseudo-période $T$ des oscillations.\\
2.4. On considère que la pseudo période $T$ est égale à la période propre $T_{0}$ de l'oscillateur $L C$.\\
Déduire la valeur de l'inductance $L$ de la bobine. (on prend ( $\pi^{2}=10$ )).

\begin{figure}[h]
\begin{center}
  \includegraphics[alt={},max width=\textwidth]{f47096eb-5ae7-4c60-ab2f-f9822564c423-4_487_864_1722_1117}
\captionsetup{labelformat=empty}
\caption{Figure 3}
\end{center}
\end{figure}

2.5. Calculer la variation de l'énergie totale $\Delta \mathscr{E}$ du circuit, entre les deux instants $t_{0}=0$ et $t_{1}=2 T$. Expliquer de point de vue énergétique le résultat obtenu.\\
2.6. Pour entretenir les oscillations électriques dans le circuit, on monte en série avec le condensateur et la bobine précédemment utilisés, un générateur ( $G$ ) qui délivre une tension $u_{G}(t)$ proportionnelle à l'intensité du courant électrique $u_{G}(t)=k . i(t)$ avec $k$ une constante (Figure 4).\\
2.6.1. Établir l'équation différentielle vérifiée par la charge $q(t)$ du condensateur.

\begin{figure}[h]
\begin{center}
  \includegraphics[alt={},max width=\textwidth]{f47096eb-5ae7-4c60-ab2f-f9822564c423-4_318_314_2316_1644}
\captionsetup{labelformat=empty}
\caption{Figure 4}
\end{center}
\end{figure}

2.6.2. Déterminer la valeur de $k$ pour que le circuit soit siège des oscillations électriques périodiques.

